\subsection{Relative contraction and the original central height}

The exact phase reduces a passage to scalar transit and normal
contraction.  Its linear normal factor may have a divergent action,
so we first estimate the nonlinear amplitude relative to that factor.
These relative estimates remain uniform as the transit time diverges.
We then compute the rates for a moving entry and connect the output
to $X=0$, where the closing difference is measured.

\begin{proposition}[Relative normal contraction]\label{rw:relative-core}
Fix $N\ge1$ and a slightly larger state and coefficient domain on which
\begin{equation}\label{rw:relative:field}
 \dot w=q f(w;\mathbf b),\qquad
 \dot Z=-Z\beta(w,Z;\mathbf b),\qquad
 \beta,\ \beta+Z\beta_Z\ge b_*>0.
\end{equation}
Here $0<q\le1$, the coefficients and the field have bounded selected
jets through order $N+2$. The triangular field has a complete auxiliary
extension preserving these bounds and the normal margins.
The selected derivatives are ordinary normal-input, slow-entry and
coefficient derivatives and, when $q=\epsilon$, Euler scale derivatives.
Suppose that their scalar flow variations through the required order
are bounded by
\begin{equation}\label{rw:relative:slow-growth}
 C_N(1+t)^{M_N}e^{c_N\varpi t}.
\end{equation}
Choose the normal neighborhood and slow width so that
$c_N\varpi<b_*/8$, increasing $c_N$ first to cover all products of
variations through this order. At the fixed slow exit suppose that its
velocity has magnitude at least $cq$ and that its hitting time satisfies
\begin{equation}\label{rw:relative:time-jets}
 |D^\alpha T|\le C_Nq^{-M_N}(1+T)^{M_N}e^{c_N\varpi T},
 \qquad 1\le|\alpha|\le N+1.
\end{equation}
If $q$ tends to zero, also require $T\ge c/q$.

For every actual passing orbit and signed normal input $z$ in a fixed
small interval, its normal output is
\begin{equation}\label{rw:relative:output}
 G(z)=Lz e^{j(z)},\qquad
 L=\exp\left[-\int_0^T\beta(w(t),0;\mathbf b)\,\dd t\right]>0,
 \qquad j(0)=0.
\end{equation}
The selected jets of $j$ through order $N$ are bounded uniformly in
$T$ and the minimum positive slow gap. The jets of $L$ and $G$ have
absolute bounds $C_Ne^{-b_*T/4}$, after increasing $C_N$; this assertion
includes mixed normal-input derivatives of $G$. In particular, normal
input zero is included in the relative statement.

If $z_*(\mathbf b)$ has bounded selected jets and stays in the smaller
normal interval, then
\begin{equation}\label{rw:relative:difference}
 G(z_*+u)-G(z_*)=G'(z_*)u e^{r(u)},\qquad
 r(0)=0,\qquad \|r\|_{N}\le C_N,
\end{equation}
on a fixed smaller $u$ interval. At a slow-root loss of the exit,
with entry and exit sections remaining separated and $T\to\infty$,
define the relative
integral along the complete auxiliary flow up to infinite time and
extend $L$ and the absolute output by zero. These definitions give
extensions of class $C^N$ in the selected ordinary variables through
each such boundary, with the stated Euler bounds.
At $\epsilon=0$ this assertion uses the Euler scale class just specified,
not ordinary scale derivatives. They do not assign a physical passage
to a nonpassing parameter.
\end{proposition}

\begin{proof}
Write $g(w,Z)=Z\beta(w,Z)$. The normal equation and its first variation
give
\[
 |Z(t;z)|\le |z|e^{-b_*t},\qquad
 G_t'(z)=\exp\left[-\int_0^t g_Z(w(s),Z(s;z))\,\dd s\right]>0.
\]
Subtracting the invariant-graph rate leaves the nonlinear defect
\begin{equation}\label{rw:relative:integral}
 j(z)=-\int_0^T\mathcal I(t,z)\,\dd t,\qquad
 \mathcal I=\beta(w(t),Z(t;z))-\beta(w(t),0)
 =Z(t;z)\int_0^1\beta_Z(w(t),sZ(t;z))\,\dd s.
\end{equation}
The mixed normal induction in Appendix~\ref{rw:app:relative-variations} gives
\begin{equation}\label{rw:relative:majorant}
 |D^\alpha\mathcal I(t,z)|
 \le C_N(1+t)^{M_N}e^{-(b_*-c_N\varpi)t},\qquad |\alpha|\le N.
\end{equation}
Its endpoint calculation uses \eqref{rw:relative:time-jets} and, for
$q\to0$, $q^{-1}\le T/c$. It bounds the relative jets and proves the
absolute estimates for $G,L$ separately.

For the difference formula, use instead
\[
 \log\frac{G'(z)}{L}
 =-\int_0^T[g_Z(w(t),Z(t;z))-\beta(w(t),0)]\,\dd t.
\]
The same majorant makes
$G'(z_*+\theta u)/G'(z_*)$ a positive unit with bounded
selected jets. Integrate it for $0\le\theta\le1$ and take its logarithm
to obtain \eqref{rw:relative:difference}.

Persistence of a finite transverse exit forces $T\to\infty$ at the
specified separated-section slow-root boundary. The appendix proves the
relative extension by differentiated tails and the absolute zero extension.
Entry--exit collision, where $T\to0$, is excluded.
\end{proof}

For the lower phase, take $q=\epsilon$ and
$\varpi=O(\epsilon)$; separated sections give $T\ge c/\epsilon$.
For an upper angular passage, take $q=1$ and use a small angular
width to obtain normal domination.  In both cases the estimates allow
independently vanishing slow gaps.  The action $-\log L$ remains in
the linear factor, with no bound required for it or its gap derivatives.

\begin{corollary}[Continuity of the relative passage]
\label{rw:transitions:relative-continuity}
Consider a family satisfying Proposition~\ref{rw:relative-core}
with common domains, margins and finite-order bounds. Suppose its
auxiliary trajectories and reserved mixed jets converge on bounded
time intervals, and its exit times converge to the corresponding
finite exit or to infinity. Then the relative amplitude and the
logarithm of the normal derivative divided by its linear factor
converge in the selected finite norms, including at a lost exit.
No derivative of the divergent linear action is included.
\end{corollary}
\begin{proof}
Use \eqref{rw:relative:majorant} and the differentiated endpoint bound
\eqref{rw:app:endpoint-term}: pass to the limit on a fixed time interval,
then send its endpoint to infinity. At a finite exit, transversality
also gives convergence of the time jets.
\end{proof}

\begin{lemma}[A moving slow entry]\label{rw:transitions:moving-rate}
In Proposition~\ref{rw:relative-core}, fix a normal input $Z_0>0$ and
the exit $v_R$. Let the slow entry $v$ vary with
$f(s)>0$ on the whole interval $[v,v_R]$. With
$\beta_0(s)=\beta(s,0)$, the actual output is
\begin{equation}\label{rw:relative:moving-output}
 G(v)=Z_0\exp\left[-q^{-1}\int_v^{v_R}
                    \frac{\beta_0(s)}{f(s)}\,\dd s+j(v)\right].
\end{equation}
The relative function $j$ has the selected jets of the preceding
proposition, with $v$ counted as an ordinary input. Put
\begin{equation}\label{rw:relative:rates}
 N_0=\beta_0+qf j',\qquad
 A_0=N_0-qf'+qf\frac{N_0'}{N_0}.
\end{equation}
Where $N_0\ge n_*>0$,
\[
 \frac{G'}G=\frac{N_0}{qf},\qquad
 \frac{G''}{G'}=\frac{A_0}{qf}.
\]
For a regular change $v=v(t)$, $v'\ne0$, the latter formula becomes
\begin{equation}\label{rw:relative:curvature-rate}
 \frac{(G\circ v)''}{(G\circ v)'}
 =\frac{A_0v'+qf\,v''/v'}{qf}.
\end{equation}
Thus bounds through order $k$ for $A_0$ use $k+2$ entry derivatives
of $j$, and the numerator after a change of input uses $k+2$
derivatives of that change and a bound for $1/v'$.
\end{lemma}

\begin{proof}
Separate the linear rate in \eqref{rw:relative:integral} and logarithmically
differentiate; the chain rule gives \eqref{rw:relative:curvature-rate}.
The denominators $N_0,v'$ require the stated margins and derivative
orders. The transit integral uses $f>0$ along the whole path.
\end{proof}

\begin{lemma}[The fixed phase exit in original coordinates]
\label{rw:transitions:curved-exit}
Choose a fixed lower phase exit away from $T=0$. Use the compact lower
coordinates and their inverse from Proposition~\ref{rw:phase}. At this
exit write $Z=Z_{\rm out}$ and set $\eta=Z_{\rm out}/\epsilon$.
On a fixed sufficiently small signed $\eta$ interval, passage from this
curved exit to a fixed original section $|X|=R$ has bounded ordinary
$\eta$, coefficient and Euler scale jets of every prescribed finite
order, provided the phase construction retains these orders. After
orienting the normal coordinate positively, the deviation from the
same weak reference on $|X|=R$ is
\begin{equation}\label{rw:relative:curved-section}
 W=\eta u_c(\eta,\mathbf b,\epsilon),\qquad
 u_c\ge c_*>0,
\end{equation}
with bounded jets of $\log u_c$. The original horizontal displacement
at the curved exit is $O_{C^N}(\epsilon\eta)$.
\end{lemma}

\begin{proof}
At the fixed phase exit, inversion gives $T=T_b+ZO_{C^N}(1)$,
with $T_b>c>0$; hence $X=-1/T$ moves by $ZO_{C^N}(1)$.
The exact compact relation is $w=\epsilon zT$ and
$w=\epsilon^2\nu T^2H(T)+Z$. Hence
\[
 z=\epsilon\nu TH(T)+\frac{Z}{\epsilon T}.
\]
Subtract the same weak reference at the shifted $T$ and substitute
$Z=\epsilon\eta$ before differentiation. The deviation is $\eta/T$,
a positive bounded unit times $\eta$.

Put $|X|=R+\epsilon\xi$. Its graph equation has bounded finite jets
on a fixed $(\xi,\eta)$ box, since $R>0$ and the horizontal denominator
has a positive margin. The section moves by $\xi=O(\eta)$; finite
ODE dependence and the positive normal derivative give
\eqref{rw:relative:curved-section} and the inverse bounds.

For an actual lower slow passage, retain a fixed final piece with
time at least $c/\epsilon$. Proposition~\ref{rw:relative-core} then
implies that $Z_{\rm out}/\epsilon$ and all selected jets are
exponentially small. Thus the actual output lies in the signed box
used for this construction.
\end{proof}

\begin{proposition}[The Gaussian central connection]
\label{rw:transitions:central-core}
Fix $0<R_0<R_1$ and a compact coefficient box such that
\[
 F_+(X)=1+\bar aX+\bar BX^2\ge f_*>0
 \quad(0\le X\le R_1).
\]
Let $R\in[R_0,R_1]$ be fixed independently of the section input, with
bounded selected coefficient and Euler jets if it varies with parameters.
Continue the chosen weak reference from
$z_b(R)=\epsilon\nu h_b$, where $h_b$ has bounded selected jets.
For sufficiently small $\epsilon>0$ and $|\nu|$, this continuation
exists on $0\le X\le R$ and satisfies
\begin{equation}\label{rw:relative:reference-weight}
 |z_b(X)|\le\frac{C\epsilon|\nu|}{X+\sqrt\epsilon},\qquad
 |D^\alpha z_b(Ry)|\le
       \frac{C_N\epsilon}{Ry+\sqrt\epsilon},\quad |\alpha|\le N.
\end{equation}
Here $0\le y\le1$ is fixed when taking coefficient or Euler
derivatives. The factor $\nu$ remains whenever no $\nu$ derivative
is taken.

For signed input $z(R)=z_b(R)+W$ in a fixed small interval, its
central height is
\begin{equation}\label{rw:relative:central-output}
 z(0)=s_b+e^{-S_c/\epsilon}a_cW e^{j_c(W)},\qquad
 s_b=z_b(0),\qquad S_c=\int_0^R\frac{X}{F_+(X)}\,\dd X>0.
\end{equation}
The positive factors $a_c,a_c^{-1}$ and the selected jets of
$\log a_c,j_c$ through order $N$ are bounded, and $j_c(0)=0$.
Ordinary $W$ derivatives are included. Composing with
\eqref{rw:relative:curved-section} gives the same statement with
$W$ replaced by $\eta$ and with adjusted positive units.
All reference choices are independent of the moving section input.
\end{proposition}

\begin{proof}
The exact original graph equation is
\begin{equation}\label{rw:relative:core-field}
 \frac{\dd z}{\dd X}=\mathfrak f(X,z)
 =\frac{(X/\epsilon)z+\bar c z(z-1)+\nu(1-z)^2}{D(X,z)},
 \quad
 D=F_++\epsilon\bar m-(1+2\epsilon\bar m)z
                         +\epsilon\bar m z^2.
\end{equation}
Its denominator stays positive on a fixed small signed tube.
The weighted Gaussian kernel and Volterra estimates in
Appendix~\ref{rw:app:relative-core} give
\eqref{rw:relative:reference-weight}, including its selected mixed jets.

Put $v=z-z_b$. Exact Taylor division gives
\[
 v_X=a(X)v+B_c(X,v)v^2,\quad
 a(X)=\frac{X}{\epsilon F_+(X)}+a_r(X),\quad
 |a_r|\le C,\quad
 |\partial_v^jB_c|\le C_j(X/\epsilon+1).
\]
These bounds include selected coefficient and Euler jets because
\eqref{rw:relative:reference-weight} controls $(X/\epsilon)z_b$.
Set $E(X)=\exp[-\int_X^R a(s)\,\dd s]$ and solve directly for the
relative logarithm:
\begin{equation}\label{rw:relative:normalized-volterra}
 Q(X)=-W\int_X^R B_c(s,WE(s)e^{Q(s)})E(s)e^{Q(s)}\,\dd s.
\end{equation}
The weighted kernel bound \eqref{rw:app:normalized-kernel-bound} and
the differentiated contraction proof in Appendix~\ref{rw:app:relative-core}
give $\|Q\|_\infty\le C|W|$ and bounded selected jets on a fixed
small $W$ interval, including zero. Thus $v(X)=WE(X)e^{Q(X)}$.
Set
$a_c=\exp[-\int_0^R a_r]$ and $j_c(W)=Q(0,W)$.
This proves \eqref{rw:relative:central-output}, retaining the exponentially
small factor in $E$ throughout.
\end{proof}

The Gaussian layer is controlled by the weighted reference estimate;
ordinary high-order $X$ derivatives there may be unbounded.  All
derivatives of the finitely smooth phase and reference are taken on
their real parameter domains.  The signed construction supplies the
estimates and coefficient comparisons, and its physical restriction
consists of positive-height original trajectories with finite
transverse passages.
For the incoming core, apply the same construction after reflection
with $(\bar a,\bar c,\nu)$ replaced by
$(-\bar a,-\bar c,-\nu)$.
